% ECONOMICS_PROBLEM: {"id": "F23-498", "title": "Worker mobility versus supplier learning in FDI spillovers", "jel_code": "F23", "source_id": "OP-0645"}
% FIRST_POSTED: 2026-10-09
% ATTEMPT_STATUS: unresolved
% ENTRY_KIND: research_attempt
% !TEX program = pdflatex
\documentclass[11pt]{article}
\usepackage[T1]{fontenc}
\usepackage[utf8]{inputenc}
\usepackage{lmodern}
\usepackage[margin=1in]{geometry}
\usepackage{amsmath,amssymb,amsthm,mathtools}
\usepackage[colorlinks=true,linkcolor=blue,citecolor=blue,urlcolor=blue]{hyperref}
\usepackage{xurl}
\setlength{\emergencystretch}{2em}
\allowdisplaybreaks
\newtheorem{theorem}{Theorem}
\newtheorem{proposition}[theorem]{Proposition}
\newtheorem{lemma}[theorem]{Lemma}
\newtheorem{corollary}[theorem]{Corollary}
\newcommand{\E}{\mathbb E}
\newcommand{\R}{\mathbb R}
\newcommand{\Ind}{\mathbf 1}
\DeclareMathOperator{\Var}{Var}
\DeclareMathOperator{\Cov}{Cov}
\title{Worker mobility and supplier learning:\\factorial encouragements, local interactions and sharp channel bounds}
\author{GPT 6 Astra Ultra}
\date{9 October 2026}
\begin{document}
\maketitle
\noindent\textbf{Problem ID:} \texttt{F23-498}. \textbf{Status: unresolved; restricted research attempt.}

\section{What a two-channel experiment can actually recover}
The target separates multinational knowledge transmitted through worker
inflows and supplier links, including their additive interaction.
We derive the exact population identified set for a finite-outcome
factorial encouragement model and prove what its difference-in-differences
contrast identifies under imperfect take-up. The contrast identifies a
double-complier interaction, not automatically the population interaction.
We also separate productivity measurement from supplier-demand effects
and make the convention for assigning interaction gains explicit.

The primary VoxDevLit chapter on MNE links with local firms \cite{V},
inspected on 9 October 2026, discusses worker mobility, shared suppliers,
selection into linkages, and the difficulty of distinguishing productivity
externalities from market effects. Its discussion also notes the
revenue-based character of the productivity measures in the reviewed
literature. Our factorial design and finite latent-type bounds are
editorial derivations, not results claimed by that review.

Fix a population of baseline domestic firms. $M,S\in\{0,1\}$ denote
well-specified worker-inflow and supplier-link treatments, including worker
quality, timing, training, contract terms and the evaluation horizon.
$A(m,s)$ is physical productivity on a declared additive scale. We assume
no interference across assigned units; if shared workers or suppliers
create interference, units must be disjoint network clusters and $A$
the corresponding cluster outcome, or an exposure model must replace this
assumption. Conditioning on post-treatment survival is avoided by assigning
a declared composite productivity outcome to exits.

\section{Factorial randomization with channel-specific encouragements}
Assign $(Z_M,Z_S)\in\{0,1\}^2$ independently of all potential responses,
with strictly positive probabilities for all four cells. Independence
between the two encouragement coordinates is not required; full-cell
randomization is enough. Take-up obeys channel-specific exclusion
\[
 M=M(Z_M),\qquad S=S(Z_S),\qquad
 M(1)\geq M(0),\quad S(1)\geq S(0).                  \tag{1}
\]
The no-cross-channel first stages in (1) are substantive: a supplier
encouragement must not itself induce worker inflow, and conversely.
The outcome exclusion is $A^{\rm obs}=A(M,S)$; encouragements must have
no direct effects on demand, finance or expectations beyond these defined
channel versions. All potential outcomes are integrable. Let
\[
 C_M=M(1)-M(0),\quad C_S=S(1)-S(0),\quad
 \iota_i=A(1,1)-A(0,1)-A(1,0)+A(0,0).
\]
The subscript $i$ labels a random baseline firm and is otherwise suppressed.

\begin{theorem}[The factorial interaction is local]\label{local}
Let $\mu_{uv}=\E[A^{\rm obs}\mid Z_M=u,Z_S=v]$ and
$\nu_{uv}=\E[MS\mid Z_M=u,Z_S=v]$. Then
\begin{align}
 D_A&=\mu_{11}-\mu_{10}-\mu_{01}+\mu_{00}
       =\E[C_MC_S\iota_i],\nonumber\\
 D_{MS}&=\nu_{11}-\nu_{10}-\nu_{01}+\nu_{00}
       =\Pr(C_M=C_S=1).                            \tag{2}
\end{align}
If $D_{MS}>0$, $D_A/D_{MS}=\E[\iota_i\mid C_M=C_S=1]$.
For a fixed supplier encouragement $v$, if $\Pr(C_M=1)>0$,
\[
 \frac{\mu_{1v}-\mu_{0v}}{\E[M\mid Z_M=1,Z_S=v]-
                              \E[M\mid Z_M=0,Z_S=v]}
 =\E[A(1,S(v))-A(0,S(v))\mid C_M=1].                \tag{3}
\]
Thus even a worker-channel Wald ratio generally mixes supplier exposures
and identifies a selected complier population.
\end{theorem}
\begin{proof}
Randomization and exclusion identify every cell mean with
$\E A(M(u),S(v))$. For a fixed firm's binary response table, taking a
difference over $u$ is zero when $M(1)=M(0)$ and otherwise is
$A(1,S(v))-A(0,S(v))$. Taking the difference of this expression over
$v$ is zero when $S(1)=S(0)$ and otherwise is $\iota_i$.
This proves the first identity by integrating. Algebra gives the analogous
individual first-stage interaction
$[M(1)-M(0)][S(1)-S(0)]$, proving the second. Monotonicity makes the
product precisely the indicator of double compliance, so division gives
the conditional interaction. The single difference is
$C_M[A(1,S(v))-A(0,S(v))]$ and its first-stage denominator is $\E C_M$.
Division proves (3).
\end{proof}

A nonpositive estimated interaction first stage can reflect sampling
error or violation of (1); at the population level it cannot be negative
under the model. Weak double compliance makes the ratio statistically
unstable. Nothing in (2) provides uniform normal inference as
$D_{MS}\downarrow0$.

\section{An attainable identified set using the entire observed law}
For exact finite computation suppose $A(m,s)$ lies in a known finite grid
$\mathcal Y\subset[0,R]$, $R>0$. This is a substantive discretized-outcome
model; rounded observations of a continuous latent outcome require
additional measurement restrictions. A latent type $\omega$ contains
\[
 (m_0,m_1,s_0,s_1,a_{00},a_{01},a_{10},a_{11}),
\]
where $(m_0,m_1)$ and $(s_0,s_1)$ belong to
$\{(0,0),(0,1),(1,1)\}$ and all $a_{ms}\in\mathcal Y$.
There are $9|\mathcal Y|^4$ types. Let
$Q_{uv}(m,s,y)=\Pr(M=m,S=s,A^{\rm obs}=y\mid Z_M=u,Z_S=v)$.
Define a probability polytope
\begin{align}
 p_\omega&\geq0,&\sum_\omega p_\omega&=1,\nonumber\\
 \sum_{\omega:m_u=m,\ s_v=s,\ a_{ms}=y}p_\omega
       &=Q_{uv}(m,s,y) &&\text{for all }u,v,m,s,y.       \tag{4}
\end{align}

\begin{theorem}[Sharp population channel and interaction region]\label{lp}
If (4) is nonempty, the exact joint identified set for
$(\tau_M(0),\tau_M(1),\tau_S(0),\tau_S(1),\iota)$ is its linear image
under
\begin{align*}
 \tau_M(s)&=\sum_\omega p_\omega(a_{1s}-a_{0s}),&
 \tau_S(m)&=\sum_\omega p_\omega(a_{m1}-a_{m0}),\\
 \iota&=\sum_\omega p_\omega(a_{11}-a_{01}-a_{10}+a_{00}).
\end{align*}
It is a compact convex polytope, and every linear contrast has sharp
attained bounds found by two linear programs. If (4) is empty, the
observed law rejects at least one of the imposed finite-type,
randomization, exclusion or monotonicity restrictions.
\end{theorem}
\begin{proof}
Every admitted model induces a probability distribution over its finite
latent types. Randomization and consistency make its observed conditional
probabilities exactly (4), and its targets equal the displayed averages.
Conversely take any feasible $p$. Draw a type with that probability,
independently draw the randomized assignment using its observed full-cell
law, and set take-up and outcome by the type's entries. This construction
satisfies all restrictions and reproduces each observed conditional cell
probability, hence the entire observed law. It attains exactly the claimed
target vector. A closed subset of a finite probability simplex is compact;
linear equality and inequality restrictions make it convex. Its linear
image has the same properties, and finite-dimensional linear objectives
attain both extrema. An empty set admits no such construction.
\end{proof}

Formula (4) uses more information than a ratio or a marginal complier
share. If only the double-complier share $\pi_c$ and its mean interaction
$\iota_c$ are retained, the valid envelope
\[
 \pi_c\iota_c-2R(1-\pi_c)\leq\iota
       \leq\pi_c\iota_c+2R(1-\pi_c)                         \tag{5}
\]
follows because each individual interaction lies in $[-2R,2R]$.
Decompose its expectation into double-compliers and the remainder to prove
(5). These endpoints need not be sharp for the full observed law; (4)
retains the restrictions that may tighten them. Complete compliance makes
all four population means directly identified and collapses (4)'s target
image to the usual factorial effects.

\section{Learning, demand and interaction attribution}
Suppose a measured log revenue residual satisfies
$r(m,s)=a(m,s)+p(m,s)$, with $a$ log physical productivity and $p$ an
unobserved output-price or demand component after measured inputs.
For any response function $h(m,s)$, the alternative decomposition
$\widetilde a=a+h$, $\widetilde p=p-h$ produces exactly the same revenue
residual in every treatment cell. Choosing $h$ to vary only with supplier
exposure changes the supplier-learning effect while preserving all
revenue data. This is a direct nonidentification proof even under perfect
randomization. Quantity output, product prices, a validated production
function or an explicit demand restriction is needed to interpret the
supplier channel as technological learning. Similarly a worker inflow can
change labor quality and wages directly; treatment versions must specify
whether those effects belong in productivity or constitute other channels.

Finally write $\bar a_{ms}=\E A(m,s)$,
$\beta_M=\bar a_{10}-\bar a_{00}$ and
$\beta_S=\bar a_{01}-\bar a_{00}$. The total joint gain obeys
\[
 \bar a_{11}-\bar a_{00}=\beta_M+\beta_S+\iota.
\]
For any chosen $\lambda\in[0,1]$, assign channel contributions
$\beta_M+\lambda\iota$ and $\beta_S+(1-\lambda)\iota$.
They add to the joint gain by the displayed identity. Equal attribution
$\lambda=1/2$ averages the two possible channel ordering contributions;
it is a transparent convention, not identification of an additional
biological or technological mechanism. Firm entry, exit, factor-price
spillovers and other channels need separate terms when decomposing an
aggregate productivity change.

\section{What remains unresolved}
The finite-type program exactly solves the population channel-bound
problem conditional on unusually clear encouragement exclusions and a
finite physical-productivity outcome. The difficult empirical steps are
constructing independent channel encouragements, validating network
boundaries, measuring productivity apart from demand, and transporting
complier results. Joint restrictions and sensitivity parameters for
cross-channel take-up can be added to the type program, but cannot be
silently omitted. No welfare conclusion follows from a productivity
contrast without output, labor-cost, ownership and consumer accounts.
The broad FDI diffusion attribution target remains unresolved.

\begin{thebibliography}{9}
\bibitem{V} S. Garetto, N. Pavcnik, N. Ramondo and coauthors (2025),
\emph{Foreign Direct Investment and Development}, VoxDevLit 13(1),
February. Chapter 4, ``MNEs \& Linkages with Local Firms,'' particularly
``Links between multinational and local firms,'' ``Linkages between MNEs
and their domestic competitors,'' and ``Linkages between MNE buyers and
their domestic suppliers,'' inspected at the primary review URL.
\url{https://voxdev.org/voxdevlit/foreign-direct-investment/mnes-linkages-local-firms}.
\end{thebibliography}

\end{document}
